\subsection{THE PRINCIPLE OF INDUCTION}

C61 (Principle of Induction). Let \(\mathbb{R}\{n\}\) be a relation in a theory \(\mathfrak{C}\) (where \(n\) is not a constant of \(\mathfrak{C}\) ). Suppose that the relation

\[
\mathrm{R} \{0 \} \text { and } (\forall n) ((n \text { is   an   integer   and } \mathrm{R} \{n \}) \Longrightarrow \mathrm{R} \{n + 1 \})
\]

is a theorem in \(\mathfrak{C}\) . Under these conditions the relation

\[
(\forall n) ((n \text {   is   an   integer   }) \Longrightarrow R \{n \})
\]

is a theorem in \(\mathfrak{C}\) .

We shall argue by contradiction. Suppose that the relation

\[
(\exists n) (n \text {   is   an   integer   and   } (\text { not   } R \{n \}))
\]

is true. Let \(q\) be an integer such that ``not \(\mathbb{R}\{q\}\) '' (method of the auxiliary constant; cf.~Chapter I, §3, no. 3 and §4, no. 1). The integers \(n\) for which `` \(n \leqslant q\) and (not \(\mathbb{R}\{n\}\) )'' form a well-ordered non-empty set (§3, no. 2, Remark), which therefore has a least element \(s\) . If \(s = 0\) , then ``not \(\mathbb{R}\{0\}\) '', contrary to hypothesis. If \(s > 0\) , then \(s = s' + 1\) , where \(s'\) is an integer such that \(s' < s\) (no. 2, Proposition 2). By definition of \(s\) , we have \(\mathbb{R}\{s'\}\) , but then the hypothesis implies that \(\mathbb{R}\{s\}\) is true, contrary to the definition of \(s\) .

In order to apply the principle of induction it is necessary in particular to prove the relation

\[
(n \text {   is   an   integer   and   } R \{n \}) \Longrightarrow R \{n + 1 \}.
\]

For this purpose the method of the auxiliary hypothesis (Chapter I, §3, no. 3) is commonly used, and it is for this reason that the relation ``n is an integer and R\{n\}'' (or even R\{n\}) is called the inductive hypothesis.

Remark. There are various criteria which are frequently used under the name of ``principle of induction''. They can all be easily deduced from C61, and we indicate here the most important of them:

\begin{enumerate}
\def\labelenumi{(\arabic{enumi})}
\tightlist
\item
  Let \(\mathbf{S}\{n\}\) be the relation
\end{enumerate}

\((\forall p)((n \text{ is an integer and } p \text{ is an integer and } p < n) \Longrightarrow \mathbb{R}\{p\})\)

and suppose that \(S\{n\}\) implies \(R\{n\}\) . Then the relation

\[
(\forall n) ((n \text {   is   an   integer   }) \Longrightarrow R \{n \})
\]

is true. For the relation \(S\{0\}\) is true, and by hypothesis \(S\{n\}\) implies \(R\{n\}\) ; since the relation \(m < n + 1\) is equivalent to \(m \leqslant n\) (no. 2, Proposition 2), the relation \(S\{n + 1\}\) is equivalent to `` \(S\{n\}\) and \(R\{n\}\) '', consequently, \(S\{n\}\) implies \(S\{n + 1\}\) . The criterion C61 therefore proves that the relation

\[
(\forall n) ((n \text {   is   an   integer   }) \Longrightarrow \mathrm{S} \{n \})
\]

is true, and since \(\mathbf{S}\{n\}\) implies \(\mathbf{R}\{n\}\) , the relation

\[
(\forall n) ((n \text {   is   an   integer   }) \Longrightarrow R \{n \})
\]

is true.

\begin{enumerate}
\def\labelenumi{(\arabic{enumi})}
\setcounter{enumi}{1}
\tightlist
\item
  Let \(k\) be an integer and let \(\mathbf{R}\{n\}\) be a relation such that the relation
\end{enumerate}

\[
\mathrm{R} \{k \} \text { and } (\forall n) ((n \text { is   an   integer } \geqslant k \text { and } \mathrm{R} \{n \}) \Longrightarrow \mathrm{R} \{n + 1 \})
\]

is true. Then the relation

\[
(\forall n) ((n \text {   is   an   integer   } \geqslant k) \Longrightarrow \mathrm{R} \{n \})
\]

is true (``induction starting at \(k\)''). For let \(S\{n\}\) be the relation

\[
(n \geqslant k) \Longrightarrow \mathrm{R} \{n \}.
\]

Then by the method of disjunction of cases we see that S\{0\} is true. On the other hand, it is easily verified that the relation

\[
(n \text {   is   an   integer   and   } S \{n \}) \Longrightarrow S \{n + 1 \}
\]

is true. It follows from C61 that the relation

\[
(n \text {   is   an   integer) } \implies \mathrm{S} \{n \}
\]

is true, which proves our assertion.

\begin{enumerate}
\def\labelenumi{(\arabic{enumi})}
\setcounter{enumi}{2}
\tightlist
\item
  Let \(a\) and \(b\) be two integers such that \(a \leqslant b\) , and let \(R\{n\}\) be a relation such that
\end{enumerate}

\[
\mathrm{R} \left\{a \right\} \text {   and   } (\forall n) ((n \text {   is   an   integer   and   } a \leqslant n <   b \text {   and   } \mathrm{R} \left\{n \right\}) \Longrightarrow \mathrm{R} \left\{n + 1 \right\}).
\]

Then the relation

\[
(\forall n) ((n \text {   is   an   integer   and   } a \leqslant n \leqslant b) \Longrightarrow \mathrm{R} \{n \})
\]

is true. The proof is similar to that of the preceding case; we take S\{n\} to be the relation ``(a ≤ n \textless{} b) ⇒ R\{n\}'' (``induction restricted to an interval'').

\begin{enumerate}
\def\labelenumi{(\arabic{enumi})}
\setcounter{enumi}{3}
\tightlist
\item
  Let \(a, b\) be two integers such that \(a \leqslant b\) , and let \(\mathbf{R}\{n\}\) be a relation such that
\end{enumerate}

\(R\{b\}\) and \((\forall n)((n \text{ is an integer and } a \leqslant n < b \text{ and } R\{n + 1\}) \Longrightarrow R\{n\}).\)

Then the relation

\[
(\forall n) ((n \text {   is   an   integer   and   } a \leqslant n \leqslant b) \Longrightarrow \mathrm{R} \{n \})
\]

is true. For we have the relation

(n is an integer and \(a \leqslant n < b\) and (not \(\mathbf{R}\{n\}\) ) \(\Longrightarrow\) not \(\mathbf{R}\{n + 1\}\) .

If for some \(n\) such that \(a \leqslant n \leqslant b\) we had (not \(\mathbf{R}\{n\}\) ), it would follow from (3) that (not \(\mathbf{R}\{b\}\) ), contrary to the hypothesis, whence the result (``descending induction'').

